\subsubsection*{3.3 Die RSA-Verschlüsselung}


Ändern nun die Blickrichtung und fragen, für welche Exponenten $e\,d$ ist
\[
  a^{e\,d} \equiv a \pmod{m}?
\]
% KORRIGIERT: Voraussetzung ggT(a, m) = 1 fehlte
Multiplizieren die Euler-Formel mit $a$ (für $\ggT(a, m) = 1$) und finden
\[
  a^{1 + \varphi(m)} \equiv a \pmod{m}
\]
also \quad $e\,d = 1$ oder $1 + \boldsymbol{\varphi}(m)$ \quad % KORRIGIERT: "=" statt "\equiv"
allgemein: \quad $e\,d \equiv 1 \pmod{\varphi(m)}$

\begin{definition}[RSA: Ver- und Entschlüsselung]
% KORRIGIERT: "=" statt "\equiv ... (mod m)" beim Ver- und Entschlüsseln
So kann man das Klartextwort $a \in \Z$ verschlüsseln: \quad $a \to c \equiv a^e \pmod{m}$\\
$c \in \Z$ ist das Codewort

entschlüsseln: \quad $a \to c \equiv a^e \to c^d \equiv a^{ed} \equiv a \pmod{m}$

Man nennt $e$ den Verschlüsselungsexponenten (encoding exponent) und\\
$d$ den Entschlüsselungsexponenten (decoding exponent).\\
Die einzige Bedingung an $e$ ist \ $e \in \Zm{\varphi(m)}^*$, also $\ggT(e, \varphi(m)) = 1$.\\
$d$ ist die Inverse von $e$ in $\Zm{\varphi(m)}^*$.
\end{definition}

\begin{beispiel}
Sei $m = 101$, $e = 13 \in \Zm{\varphi(m)}^*$, \quad $\varphi(101) = 100$.\\
$13 \cdot 7 \cdot 11 = 1001 \equiv 1 \pmod{100}$, also\\
$d = 77$. Verschlüsselt werden soll die Zahl $a = 8 \in \Zm{m}^*$
\[
  c \equiv a^e = 8^{13} \equiv 18 \pmod{101} \qquad (\text{Codewort})
\]
\begin{align*}
  & 13 = 1 + 2^2 + 2^3, \\
  & 8^{13} = 8^8 \cdot 8^4 \cdot 8 \\
  & 8^2 = 64 \\
  & 8^4 = 4096 \equiv 56 \pmod{101} \\
% KORRIGIERT: Original "8^8 = 3136"; richtig ist 8^8 \equiv 56^2 = 3136
  & 8^8 \equiv 56^2 = 3136 \equiv 5 \pmod{101} \\
  \Rightarrow\quad & c \equiv 8 \cdot 56 \cdot 5 = 2240 \equiv 18 \pmod{101} && (\text{Codewort}) \\
% KORRIGIERT: "=" statt "\equiv"; Begründung ergänzt
  & a \equiv 18^{77} \equiv 8^{13 \cdot 77} \equiv 8 \pmod{101} && (\text{Klartextwort; } 13 \cdot 77 \equiv 1 \ (\bmod\ 100),\ \ggT(8, 101) = 1)
\end{align*}
\end{beispiel}

\subsubsection*{Alice, Bob und Oscar}

% KORRIGIERT: "=" statt "\equiv ... (mod m)"
Alice verschlüsselt $a \to c \equiv a^e \pmod{m}$ und sendet $c$ an Bob. \hfill (Kundin einer Bank)\\
Bob ist im Besitz von $d$ und entschlüsselt $c \to a \equiv c^d \pmod{m}$ \hfill (Bank)\\
Oscar versucht den Klartext $a$ zu bekommen.

\begin{beispiel}
Hätte Oscar eine Chance an $a$ zu kommen, wenn er sich $m$ und das verschlüsselte Wort $c$ besorgt?

Ja, ganz einfach. Oscar stellt fest, dass $m = 101$ eine Primzahl ist, dann ist $\varphi(m) = m - 1$. Oscar rechnet selbst $d$ aus
\[
% KORRIGIERT: "a = c^d" statt "a \equiv c^d (mod m)"
  e\,d \equiv 1 \pmod{\varphi(m)} \qquad \Rightarrow \qquad a \equiv c^d \pmod{m}
\]
\end{beispiel}

Wie kann das verhindert (oder sehr schwer gemacht) werden?

\begin{definition}[Öffentlicher und privater Schlüssel]
\textbf{Öffentlicher Schlüssel}\\
Alice muss verschlüsseln können
\[
% KORRIGIERT: "=" statt "\equiv"
  c \equiv a^e \pmod{m}
\]
$\Rightarrow$ \quad öffentlich ist $e$, $m$

\textbf{Privater Schlüssel}\\
Bob muss entschlüsseln können
\[
% KORRIGIERT: "=" statt "\equiv"
  a \equiv c^d \pmod{m}
\]
Bob hat $d$ aus
\[
  e\,d \equiv 1 \pmod{\varphi(m)}
\]
berechnet.\\
$\Rightarrow$ \quad privat ist $d$, $\varphi(m)$
\end{definition}

Wie macht es Bob dem Oscar schwer $\varphi(m)$ zu berechnen? Bob setzt
\begin{align*}
% KORRIGIERT: Voraussetzung p \neq q fehlte
  & m = p \cdot q, && p, q \text{ große Primzahlen, } p \neq q. \\
  \Rightarrow\quad & \varphi(m) = (p - 1)(q - 1)
\end{align*}
Aber $p$, $q$ und damit auch $\varphi(m)$ kennt nur Bob.

\begin{beispiel}
Bob wählt $p = 101$, $q = 113$\\
$m = 11413$, \quad $\varphi(m) = 100 \cdot 112 = 11200 = 2^6 \cdot 5^2 \cdot 7$.\\
$e \in \Zm{\varphi(m)}^*$, darf also nicht durch 2, 5 oder 7 teilbar sein. Setzen $e = 3533$.\\
In der Praxis prüft man $e \in \Zm{\varphi(m)}^*$ bei der Berechnung von $d$ mit dem EA.
\begin{align*}
  & d \equiv e^{-1} \pmod{\varphi(m)} \qquad (\text{existiert, da } \ggT(e, \varphi(m)) = 1) \\
  \Rightarrow\quad & d = 6597
\end{align*}
Alice wählt die Zahl $a = 9726 \in \Zm{m}^*$ und berechnet.
\begin{align*}
% KORRIGIERT: Original "c \equiv a^e = 9726^{3533} \equiv (mod m) = 5761" und "a = 5761^{6597} (mod m) = 9726"
  c &\equiv a^e = 9726^{3533} \\
    &\equiv 5761 \pmod{m}
\end{align*}
% KORRIGIERT: Original "Alice sendet c an Bob und rechnet"; entschlüsseln kann nur Bob (er kennt d)
Alice sendet $c$ an Bob. Bob rechnet
\begin{align*}
  a &\equiv 5761^{6597} \\
    &\equiv 9726 \pmod{m}
\end{align*}
\end{beispiel}

\begin{bemerkung}[Fazit]
Solange es unmöglich ist $m = p \cdot q$ zu faktorisieren, ist RSA ein sicheres Verschlüsselungsverfahren.
\end{bemerkung}
